\subsection{Agent-based economics \label{sec:ace}}

This thesis uses agents to look for spatial economic outcomes from actors with heterogenous locations. There is very little in the way of `complex' behaviour, and where it does exist (for instance in the interaction of price-setting and production), it is not the focus of analysis. This is not to deny the uses of \gls{cas} theory. Evolutionary ideas are particularly important for developing an understanding of how economic growth actually functions (\citealt{martin_complexity_2007}, and see section \ref{sec:coord_production}). The evolution of diverse market structures themselves come under the microscope (see e.g. \citealt{mirowski_markets_2007} discussing `marketomata'). But a perceived need to keep \gls{abm} and complexity umbilically linked cuts off some of the more mundane  - yet crucial - uses of heterogeneity. 

In Otter \emph{et al}'s agent paper presenting an economic geography model, for instance, it is suggested that ``regional economics and geography'' are apparently ``not sufficient to explain the complex spatial patterns, such as clusters and sprawl, that we encounter'' (\citealt{otter_abloom:_2001} p.1). They argued it is possible to use ``complexity as underlying explanatory variable'' (ibid). This generally translates into models that claim to be `validated' if they manifest some element considered complex. (Not a surprising conclusion given some of the theories of explanation that have grown along with \gls{cas} theory; see section \ref{sec:generative} on Epstein.) The `complexity' cart seems to be pulling the research horse in some cases.  Boschma and Martin, for example, ask, ``in what sense can complexity theory notions (or metaphors), such as the emergence, self-organisation, criticality and so on, be used to conceptualize the economic landscape?'' (\citealt{boschma_editorial:_2007} p.541)

As O'sullivan and Haklay say, \gls{ace} is driven by ``a view of the `economy as an evolving complex system' promoted by the Santa Fe Institute'', accompanied by a ``widespread disillusion with neoclassical equilibrium economics'' \citealt{osullivan_agent-based_2000} p.1412. Disillusion is a gentle term: the reality for some researchers is a rather more `year zero' feel, seeing \gls{cas} as ``a pioneering break from a moribund Newtonian worldview'' (\citealt{manson_simplifying_2001} p.412). Mirowski, for example, argues (paraphrased by Blaug) that ``the whole of neo-classical economics ever since the \gls{margrev} has been an attempt to create an economics that emulates all the essential features of of nineteenth-century physics.'' (\citealt{blaug_economic_1997} p.284) 

As Blaug notes, this is rather strong; there was no conscious attempt to emulate physics. Nevertheless, from a \gls{cas} perspective, the faults with `traditional' economics are seen to be self-evident. What Conlisk calls the ``strange sacrifices required for the `ritual purity' of optimisation-only models'' (\citealt{conlisk_why_1996} p.686) are no longer considered necessary. While understandable at a time when no computers were available, the feeling is - as Mirowski says - economists should ``kick the habit of their physics envy and join the 21st century by rethinking the importance of computation and evolution in the way that they approach markets'' (\citealt{mirowski_kicking_2007} p.359). Edmonds goes even further, suggesting that numerical representation itself is questionable. For example, he criticises numerical representation of variety on the basis that not all the dynamics associated with variety (such as evolutionary dynamics) can be described numerically
%``what starts out as a working model or `proxy' for the phenomena ends up being mistaken for the truth about that phenomena'' (\citealt{edmonds_against_2004} p.5).

\cmt{need section on GE, and to mention its roots and a siren to call economists to geography...}
As section x discussed, \gls{ge} shares this `Newtonian' lineage. Unsurprisingly, then, there is very little overlap between the research agendas of \gls{ge} and \gls{ace}. \gls{ace} as a distinct sub-discipline of \gls{abm} takes it name from the broader school of computational economics. Despite arguments that \gls{ace} is perfect for re-examining the Marshallian roots of economics (\citealt{leijonhufvud_agent-based_2006}), it has tended to follow the `Santa Fe' research agenda. One editorial contains a good checklist of the concepts of this agenda: ``complexity, evolution, auto-organisation... emergence... bounded rationality, inductive reasoning'' (\citealt{consiglio_artificial_2007} p.vi).  

The economic coordination problem in particular lends itself to \gls{abm}'s default focus on interacting agents. Indeed, Tesfatsion defines \gls{ace} as ``the computational study of economies modelled as dynamic systems of interacting agents''. The natural question for \gls{ace} theorists thinking about markets is then: ``who does the job of the so-called market adjustment?'' (\citealt{consiglio_artificial_2007-1} p.102). This requires the modeller to ``analyze explicitly how agents interact with each other'' (\citealt{kirman_evolving_2001} p.460) Howitt, for one, sees this requirement as a key strength of the approach: ``one of the virtues of the \gls{ace} approach to economics... is that it forces one to make explicit the mechanisms through which individual actions are coordinated, for better or worse'' (\citealt{howitt_coordination_2006} pp.1068). Agents, it is hoped, force theorists to crack open the problem and look at the dynamics inside.
%underneath, and in doing so either validate or challenge equilibrium analyses. 

Hayek's short essay, `the use of knowledge in society' (\citealt{hayek_use_1945} p.529), is something of a talisman for this way of thinking. In particular, his strong aversion to macro-level assumptions fits \gls{ace} perfectly. He insisted that ``we must show how a solution is produced by the interactions of people each of whom possesses only partial knowledge'' (\citealt{hayek_use_1945} p.530) - articulating both an assumption that the `solution' arises from micro-level interactions, and that it must come from `boundedly rational' actors. Epstein's `generativist's question' (discussed in-depth in section \ref{sec:generative}) -  ``how could the decentralised local interactions of heterogeneous autonomous agents generate the given regularity?'' (\citealt{epstein_generative_2006} p.5) - echoes Hayek's take. Some theorists (\citealt{vriend_was_2002} p.2; \citealt{miller_complex_2007}) can imagine that, had Hayek only been able to access present technology, he would surely have embraced it, and complexity theory along with it.

This Austrian tilt to \gls{ace} is further reason for the lack of overlap with the kinds of questions \gls{ge} asks. A good example of how this manifests itself is the criticism levelled at the concept of equilibrium. As Blaug notes, ``the desideratum of any economic theory is the delineation of an equilibrium end-state'' (\citealt{blaug_trade-off_2009} p.224). These end-states are analysed in the same way that basic physical models of stationary systems assume all forces balance to zero and are then able to deduce force values required for that stationary state. In-keeping with the attacks on its Newtonian underpinning, Colander argues:

\begin{quote}
``the self-correcting `stability' vision cultivated by economic pedagogy is problematic in several respects. First and foremost, it is simply wrong: stability is not the norm in complex systems'' (\citealt{colander_sins_2010} p.286).
\end{quote}

For example, the idea of the `Walrasian auctioneer' gets set up often as a target. It is described as an actual mechanism: a centralised adjudicator capable of mediating a process of `groping' towards a set of market-clearing prices, and knowing the moment when they are all correct\footnote{See e.g. \citealt{tesfatsion_agent-based_2006} p.834, where it is argued that ``equilibrium values... are determined by market clearing conditions imposed through the Walrasian Auctioneer pricing mechanism; they are not determined by actions of consumers, firms, or any other agency supposed to actually reside within the economy''; (\citealt{consiglio_artificial_2007-1}.}. Though it does not originate in Walras' own work (\citealt{blaug_economic_1997} p.555-6), the auctioneer concept is a good summary of everything \gls{ace} opposes as implausible. Once such convenient (but, it is argued, impossible) assumptions are shown to be unjustifiable, a Pandora's box is opened:

\begin{quote}
``The modeller must now come to grips with challenging issues such as asymmetric information, strategic interaction, expectation formation on the basis of limited information, mutual learning, social norms, transaction costs, externalities, market power, predation, collusion, and the possibility of coordination failure.'' (\citealt{tesfatsion_agent-based_2006} p.836)
\end{quote}

Each of these `challenging issues' does in fact garner plenty of attention from analytic economics: Stiglitz got his Nobel prize for his work on asymmetric information going back to the seventies (\citealt{rothschild_equilibrium_1976}), game theory deals with strategy (and `predation and collusion' might be considered forms of strategy), and Krugman himself discusses expectation and the role of history  as a key element of his geographical work (\citealt{krugman_history_1991}). As regards agglomeration economies, the importance of externalities has been discussed in chapter \ref{chapter:location}. 

\gls{ace} models often share one vital feature with the equilibrium methods they reject, however: a market clearing point. This often means that auction-style models are used where markets reach a defined end-point via a bidding process - through many interactions, perhaps, but still terminating at the end of trade. An example would be Kirman and Vriend's fish market model, with an interdependent morning and afternoon's trading (\citealt{kirman_evolving_2001} pp.467) The question of what happens if markets do not `clear' was asked by Hicks, and Isard understood its implications for trade across space. If all actors are making market decisions on the head of a pin, an auction-style clearing market presents itself as the natural way to think about it. But as soon as actors have heterogenous location, that becomes impossible: `clearing market' conditions do not hold. As Leijonhufvud puts it, quoting Hicks -

\begin{quote}
``In a normal, ongoing market, transactors are not all brought together in a single location and at the same time. Without centralisation and synchronisation, the supply-equals-demand condition `cannot be used to determine price, in Walras' or Marshall's manner'.'' (\citealt{leijonhufvud_agent-based_2006} p.1633; see also \citealt{ladley_effects_2007} p.83-4) 
\end{quote}

While easier to ignore this problem if geography is not an element of the model, one cannot consider spatially varied trade as `one market' (\citealt{isard_location_1956} p.43). From an \gls{ace} point of view, Dibble (\citeyear{dibble_computational_2006} p.1516), notes these effects of space on standard market-clearing assumptions. Ladley and Bullock also point out (\citealt{ladley_strategic_2008} p.296) that market actors may be segregated by space, able only to interact with a given subset of others, and that spatial segregation forms natural networks as actors' ranges of interaction overlap. (Indeed, one model does the reverse, and uses actors with randomised location in an $R^2$ space with a fixed radii as a proxy for building a social network distance model; \citealt{hamill_social_2009}, \citeyear{hamill_simulating_2010}). 

It is the `non-clearing' part of this problem that the thesis model attempts to deal with, rather than the network element. Section \ref{sec:routedev} already made the argument that transport networks can be replaced with a Euclidean proxy when discussing distance, though only by acknowledging the severe limitations for discussing change this imposes. Ladley and Bullock note Wilhite making a related point: that if network change is slow compared to activity on the network, its structure becomes more important, not less (\citealt{ladley_strategic_2008} p.299 referring to \citealt{wilhite_economic_2006}). This thesis cannot answer this issue fully as it works within the Euclidean assumptions from section \ref{sec:routedev}, where it was argued that the success of Euclidean stand-ins is dependent on both the spatial and temporal scale under examination, as well as whether network dynamics are part of the research question. One may be asking, as Epstein does, how ``the endogenous connectivity - the topology - of a social network affect its performance as a distributed computational device, one that... computes price equilibria, or converges to (computes) social norms, or converges to spatial settlement patterns such as cities?'' (\citealt{epstein_generative_2006} p.17-18). 

\subsection{Spatial versus economic agents}

A Venn diagram \gls{ace} and spatial \gls{abm} would perhaps reveal just as little overlap as between \gls{ace} and \gls{ge}.  Torrens' recent review of `\gls{abm} in the Spatial Sciences' (\citeyear{torrens_agent-based_2010}) illustrates the point. One economic example picked up on is not actually about spatial modelling: Torrens cites Farmer and Foley's Nature piece who argue that the economy is complex, \gls{abm} can `do' complexity where equilibrium models cannot, representing something closer to reality; therefore, \gls{abm} is better:

\begin{quote}
``Agent-based models potentially present a way to model the financial economy as a complex system, as Keynes attempted to do, while taking human adaptation and learning into account, as Lucas advocated. Such models allow for the creation of a kind of virtual universe, in which many players can act in complex - and realistic - ways'' (\citealt{farmer_economy_2009} p.685-6).
\end{quote}

They suggest an effort akin to general climate modelling for the economy is in order, coupling the various social science disciplines in the same way climate models might couple ocean and atmosphere feedback. Their emphasis on  building a `virtual universe' can be seen in many of the spatial models Torrens surveys (this idea is the subject of section \ref{sec:virtualworlds}). The cellular automata models surveyed by Torrens in particular have this quality. Crowd and `swarm' models loom large, leveraging the interactive properties of agents (a point noticed by O'Sullivan and Haklay, \citeyear{osullivan_agent-based_2000} p.1411). 

In terms of spatial \emph{economic} models, Torrens cites his own work (among others)  - in one example, a technically rich implementation of a combined cellular-automata agent hybrid model (\citealt{torrens_geographic_2005}). There appears to be no driving research agenda, however, beyond proving the technical feasibility of doing so, though the ability of disaggregated models to produce `emergence and self-organisation' is mentioned (p.396).
%\footnote{He also mentions Krugman's `self-organizing economy' as containing an agent-based analysis - it doesn't.}

Agent models using explicitly spatial economic ideas that would seem familiar to a \gls{ge} theorist are few. Sasaki and Box  use \gls{abm} to attempt to `verify' classic spatial analytics like the von Thunen model (\citealt{sasaki_agent-based_2003}) by showing how the spatial equilibrium result can come about through individual-level action but, aside from referring to von Thunen himself, it sticks to the \gls{abm} literature. The work of Christopher Fowler stands out as a theoretically grounded agent model with the explicit aim of examining the \gls{coremodel}; this is discussed in-depth in section \ref{sec:fowler}.
